\documentclass[12pt]{exam} 
\renewcommand{\partshook}{
	\setlength{\itemsep}{5mm}
}
\pointsinmargin
\addpoints

\usepackage{cmbright}

\pagestyle{empty} %for handouts
\addtolength{\jot}{2em} % this makes all formulae a bit more spaced vertically for handouts
\setlength{\parindent}{0pt} %we're not writing a novel
%\renewcommand{\arraystretch}{2} %makes tables taller, so easier to write in
\usepackage[a4paper, margin=1in]{geometry} %makes it easy to specify where to put stuff on the page

\usepackage[fleqn]{amsmath}
\usepackage{amssymb} %standard classy maths symbols

\usepackage[utf8]{inputenc}

\usepackage{tcolorbox}
\usepackage{systeme}

\begin{document}

\textbf{Name:} \dotfill

\begin{questions}


\question[6] Solve the following simultaneous equations

\fullwidth{

\begin{parts}
\begin{minipage}{0.45\textwidth}
\part \hspace{2mm} \systeme{ 1.0x + 1.1y = 2.9, 1.0x + 0.9y = 3.1}
\fillwithdottedlines{78mm}
\end{minipage}
\hspace{2mm}
\vrule
\hspace{8mm}
\begin{minipage}{0.45\textwidth}
\part \hspace{2mm} \systeme{ 1.0x + 1.1y = 2.9, 0.9x+1.0y = 3.0}
\fillwithdottedlines{78mm}
\end{minipage}
\end{parts}
}

\

\question[3] Make $y$ the subject \emph{(isoler $y$)} in the following equations :
\vspace{2mm}

$1.0x + 1.1y = 2.9 \Rightarrow$ \dotfill

$1.0x + 0.9y = 3.1 \Rightarrow$ \dotfill 

$0.9x + 1.0y = 3.0 \Rightarrow$ \dotfill 

\question[1] The equations in question 1 are very similar. Why are the solutions so different?
\fillwithdottedlines{\stretch{1}}


\question[0] \textbf{Bonus.} Solve $x(x^2-1) = 0$. \hfill \emph{Work on the back.}

\end{questions}


\begin{tcolorbox}

\textbf{Name of marker:}

\begin{center}
\gradetable[h][questions]
\end{center}

\end{tcolorbox}


\end{document}